\section{Framework: one estimand, four shocks, two channels, six rulebooks}\label{sec:framework}

\subsection{The estimand}

For every episode the object is the same conditional quantity: the
statutory household response to a declared trade-shock first stage. Let a
first stage specify (i) a vector of consumer-price changes by consumption
category, and/or (ii) a labour-income loss or gain with a sectoral
location and an adjustment margin. The second stage maps that first stage
onto households and computes the change in disposable income, its
distribution, and the Exchequer's share, under the tax--benefit rules in
force at the episode date. The headline metric is the \emph{cushioning
rate}---one minus the household disposable-income change per pound of
gross shock, the income-stabilisation coefficient of \citet{dolls2012}
generalised to price shocks by measuring the gross shock as the
expenditure-weighted cost imposed at fixed consumption---because ratios,
unlike levels, are comparable across episodes whose gross sizes span two
orders of magnitude.

Nothing in the design estimates the first stages. Each is imported from
published estimates and labelled with its source and its attribution
caveats (Section~\ref{sec:episodes}); the paper's contribution begins
where those estimates stop, at the household and the Exchequer. This is
the conditional-estimand convention of the companion tariff study
\citep{ahmadi2026}, applied retrospectively---with the advantage that
retrospective first stages can be anchored to realised outturns rather
than ex-ante bridges.

\subsection{Two channels}

Trade shocks reach households through prices and through earnings, and
the post-Covid episodes split cleanly on that line. The consumption
channel prices a category-level price vector against household
expenditure shares: implemented at the household level on imputed
consumption (the microsimulation stage) and, in this paper's first-pass
results, on published expenditure-by-decile tables. The earnings channel
imposes sectoral labour-income changes on employees by industry and runs
them through the statutory system, with the adjustment margin (wage cuts
versus displacement) as a declared scenario dimension---the machinery,
conventions and caveats of the companion study carry over unchanged. The
channels are reported separately and never summed without an explicit
interaction note: a price shock erodes real income that the benefit
system partially indexes, while an earnings shock triggers means-tested
entitlement; the two responses run through different statutory
mechanisms, and their comparison is the point.

\subsection{The rulebook axis}

The same shock cushions differently depending on when it lands, because
the rulebook moves: benefit uprating follows lagged inflation, so a
fast price shock outruns indexation for a year; thresholds frozen in
cash terms tighten in real terms through an inflation episode; and
discretionary interventions (most dramatically the Energy Price
Guarantee) rewrite the effective schedule mid-shock. The design
therefore separates three quantities for each episode: the cushioning
delivered by the rules \emph{as they stood}, the cushioning the same
rules would have delivered under contemporaneous full indexation (the
uprating-lag cost), and the discretionary layer on top. No study of
trade shocks anywhere has, to our knowledge, held the shock fixed and
varied the rulebook; the post-Covid sequence---four shocks landing on
four different vintages of the same welfare state---is a natural
laboratory for exactly that comparison.

\subsection{What the design is not}

The estimates are first-round statutory incidence, not general
equilibrium: no behavioural response, no wage or employment feedback
from the price shocks, no substitution in consumption (the price-channel
numbers are first-order, envelope-theorem approximations to welfare
changes and upper bounds on avoidable cost). Episode first stages carry
their own literatures' uncertainty, which the design brackets by
importing low/central/high variants where the literature provides them
and labelling single-point stages as such. The comparison across
episodes is a comparison of conditional scenarios that share a second
stage, not a causal decomposition of realised household outcomes over
2021--2026, which would require separating the trade shocks from
everything else that happened---a task this design deliberately
declines.
